\documentclass[11pt]{amsart}
\usepackage{amsmath,amssymb,mathtools,booktabs,microtype}
\usepackage[hidelinks]{hyperref}
\usepackage{xurl}
\emergencystretch=3em
\newtheorem{theorem}{Theorem}[section]
\newtheorem{lemma}[theorem]{Lemma}
\newtheorem{proposition}[theorem]{Proposition}
\theoremstyle{definition}
\newtheorem{definition}[theorem]{Definition}
\newtheorem{remark}[theorem]{Remark}
\numberwithin{equation}{section}
\newcommand{\K}{\mathbb K}
\newcommand{\Q}{\mathbb Q}
\newcommand{\Sym}{\operatorname{Sym}}
\newcommand{\Span}{\operatorname{span}}
\newcommand{\im}{\operatorname{im}}
\newcommand{\rank}{\operatorname{rank}}
\newcommand{\perm}{\operatorname{perm}}
\newcommand{\CR}{\operatorname{ChowRank}}
\newcommand{\cD}{\mathcal D}
\title{The Chow Rank of the Five-by-Five Permanent}
\author{Rongyu Xu}
\date{October 16, 2026}
\subjclass[2020]{Primary 15A69; Secondary 14N05, 68V05}
\keywords{permanent, Chow rank, product rank, catalecticant}
\hypersetup{pdftitle={The Chow Rank of the Five-by-Five Permanent},pdfauthor={Rongyu Xu}}
\begin{document}
\begin{abstract}
We prove that the ordinary Chow rank of the permanent of a generic five by
five matrix is sixteen over every field of characteristic zero. The lower
bound uses an image-span inequality for rectangular catalecticants and a
quotient-symbol estimate for a single quintic product. All mathematical
steps are proved by coefficient identities, elementary specialization, and
rank inequalities. No finite enumeration or computer-generated matrix-rank
certificate is a premise. The theorem makes no assertion about border Chow
rank.
\end{abstract}
\maketitle

\section{Statement and derivative spaces}
Let \(\K\) initially be algebraically closed of characteristic zero, and put
\[
 V=\Span_{\K}\{x_{ij}:1\le i,j\le5\},\qquad
 f=\perm_5=\sum_{\sigma\in S_5}\prod_{i=1}^5x_{i,\sigma(i)}.
\]
Here \(S_5\) is the permutation group on \(\{1,\ldots,5\}\).
For a degree-\(m\) form \(g\in\Sym^m V\), define
\[
 \begin{aligned}
 C_{a,m-a}(g)&:\Sym^a V^*\longrightarrow\Sym^{m-a}V,\\
 C_{a,m-a}(g)(\lambda_1\cdots\lambda_a)
   &:=\partial_{\lambda_1}\cdots\partial_{\lambda_a}g,\\
 \cD_j(g)&:=\im C_{m-j,j}(g).
 \end{aligned}
\]
The variables are members of \(V\), hence linear functions on \(V^*\).
For a subspace \(G\) of forms, \(\partial G\) is the span of all its
first directional derivatives. For a quadratic form, rank means the rank of
its symmetric bilinear matrix. Its essential-variable space is the least
\(L\subset V\) supporting it; this has dimension equal to its matrix rank.

The ordinary Chow rank \(\CR_{\K}(g)\) is the least number of nonzero
products of \(m\) linear forms summing to \(g\), with \(\CR_{\K}(0)=0\).
It is also called product rank \cite{IltenTeitler2016}.
Han and Song \cite{HanSong2026} prove rank sixteen for the five by five
permanent tensor, regarded as a tensor with five separate row factors.
Tensor decompositions require each linear factor to belong to its assigned
row space; an ordinary Chow decomposition imposes no such requirement.
Thus that tensor-rank lower bound does not by itself give the Chow-rank
lower bound proved here.
\begin{theorem}\label{thm:main}
Over every field \(\K\) of characteristic zero,
\[\CR_{\K}(\perm_5)=16.\]
\end{theorem}

Put \(E_j=\cD_j(f)\). We will use the following elementary facts.
\begin{lemma}\label{lem:permanent}
One has \(\dim E_2=\dim E_3=100\), \(\partial E_2=V\), and
\[
 E_2^{(1)}:=\{g\in\Sym^3 V:\partial g\subset E_2\}=E_3.
\]
Every nonzero quadratic in \(E_2\) has matrix rank at least four. If
\(L\subset V\) has dimension at most five, then
\[\dim(E_2\cap\Sym^2L)\le1.\]
\end{lemma}
\begin{proof}
The subpermanents on two rows and two columns form a basis of \(E_2\),
and those on three rows and three columns form a basis of \(E_3\).
Their matching supports are disjoint, and both counts are \(\binom52^2=100\).
Differentiating a rectangle at one corner gives the opposite variable.

If every derivative of a cubic lies in \(E_2\), a monomial with a repeated
row or column is impossible: some first derivative would contain a monomial
outside the matching support of \(E_2\). Inside one three by three block,
the rectangle relations equate the coefficients of matchings differing by a
transposition. Transpositions connect the six matchings. Thus the cubic is
a linear combination of the three by three subpermanents, proving the
prolongation assertion.

The row-column diagonal torus acts with distinct characters on the rectangle
subpermanents. A generic one-parameter subgroup takes any nonzero quadratic
in \(E_2\), projectively, to a single rectangle. Matrix rank cannot increase
under this specialization, and a rectangle has rank four.

For completeness, the coordinate-specialization argument needed for the
intersection bound is as follows. Choose a one-parameter subgroup whose
weights on the 25 variables are distinct. In the Grassmannian of
\(a\)-planes, the limit of any \(L\) is a coordinate \(a\)-plane: the
limit is fixed by the subgroup, whose variable eigenspaces are one-dimensional.
The dimension of \(E_2\cap\Sym^2L\) is upper semicontinuous, since it is
a kernel dimension of a vector-bundle map. It suffices to bound it at the
coordinate limit. A coordinate plane corresponds to a bipartite graph with
\(a\le5\) edges, and its intersection dimension counts four-cycles.
Two different four-cycles use at least six edges. The bound is therefore one.
\end{proof}

\section{A quotient-symbol estimate for quintic products}
Let \(T=\ell_1\cdots\ell_5\ne0\), and set
\[
 \begin{gathered}
 L=\Span\{\ell_1,\ldots,\ell_5\},\qquad
 F=\cD_2(T),\qquad U=\cD_3(T),\\
 u=\dim U=\dim F\le10.
 \end{gathered}
\]
The equality of dimensions follows by transposing the two catalecticants
under the perfect differential pairing in characteristic zero. The upper
bound follows because third derivatives are spanned by the ten pair products.
Let \(R=F\cap E_2\). Lemma~\ref{lem:permanent} gives \(\dim R\le1\),
and each nonzero element of \(R\) has quadratic rank at least four.

For a surjection \(P:L\twoheadrightarrow D\), of rank \(d\), let
\(q_R:F\to F/R\) be the quotient map. Polarization is
\[
 \delta:U\longrightarrow L\otimes F,\qquad
 g\longmapsto\sum_j y_j\otimes\partial_{y_j^*}g,
\]
where \((y_j)\) is any basis of \(L\) and \((y_j^*)\) its dual.
It is basis-independent. Write \(\beta=(P\otimes q_R)\delta\).

\begin{lemma}[Small complementary span]\label{lem:kernel}
If \(k=\dim\ker P\le2\), then
\[
 \ker\beta=U\cap\Sym^3(\ker P),\qquad
 \rank\beta\ge\max\{0,u-\tbinom{k+2}{3}\}.
\]
\end{lemma}
\begin{proof}
Choose coordinates consisting of a basis of \(S=\ker P\) and \(d\)
complementary variables. For a kernel cubic, its derivatives in all the
complementary directions lie in \(R\). If one is nonzero, changing the
complementary coordinates makes these derivatives \(q,0,\ldots,0\),
with \(0\ne q\in R\). Characteristic zero then implies that the cubic,
and hence \(q\), uses only \(S\) and the first complementary variable.
This gives \(\rank q\le k+1\le3\), a contradiction. All complementary
derivatives vanish, which is precisely membership in \(\Sym^3S\).
The dimension bound follows.
\end{proof}

\begin{lemma}[Five independent factors]\label{lem:independent}
If \(\dim L=5\), then \(u=10\) and \(\rank\beta\ge2d\).
\end{lemma}
\begin{proof}
Use the factors as coordinates \(z_1,\ldots,z_5\). Then \(U\) and \(F\)
are the squarefree cubic and quadratic spaces. For arbitrary
\((\ker P,R)\), simultaneously specialize in their Grassmannians by a
generic one-parameter subgroup of the diagonal torus. Its weights on the
variables and on the quadratic monomials can be chosen distinct, so both
limits are coordinate subspaces. The rank of the induced map of the
tautological quotient bundles cannot increase in this limit.
No rank-four condition on the limiting \(R\) is required here.

At a coordinate point let \(B\subset[5]\) index the kernel of \(P\),
so \(|B|=5-d\). The limiting subspace \(R\) is either zero or a
coordinate edge \(A=\Span(z_az_b)\). The source of the polarization map
has a basis indexed by the ten triples \(I\subset[5]\). The column indexed
by \(I\) has one component for each \(i\in I\), in coordinate
\(z_i\otimes z_{I\setminus\{i\}}\), and is killed precisely when every
component is removed by \(P\) or by quotienting out \(A\). Thus a column
is zero exactly when either \(I\subset B\), or \(|I\setminus B|=1\)
and the complementary edge is \(A\). The second case is possible only
when \(A\subset B\), and then there are exactly \(d\) such triples.
Distinct triples have disjoint polarization supports, so all nonzero
columns are independent. Therefore
\[
 \rank\beta\ge 10-\binom{5-d}{3}
       -d\,\mathbf1_{\{A\subset B\}},
\]
where the indicator is zero if \(R=0\). Minimizing over the coordinate
possibilities gives
\[
\begin{array}{c|rrrrrr}
d&0&1&2&3&4&5\\\hline
\rank\beta\text{ lower bound}&0&5&7&7&10&10.
\end{array}
\]
These values are respectively at least \(2d\). This proves the required
inequality directly; no 352-case enumeration is needed.
\end{proof}

\begin{lemma}[Four-dimensional factor span]\label{lem:four}
If \(\dim L=4\), then \(u\ge7\). For \(d=1\),
\[
 \rank\beta+(10-u)/2\ge2.
\]
\end{lemma}
\begin{proof}
Choosing four independent factors and rescaling coordinates puts the product
in one of the forms
\[
 T_s=x_1x_2x_3x_4(x_1+\cdots+x_s),\qquad 1\le s\le4.
\]
We determine the ranks of the ten second derivatives by coefficient
comparison, not by matrix elimination. Write \(H_{ij}=\partial_i\partial_jT_s\).
For \(s=1\), the nonzero derivatives are
\[
\begin{gathered}
H_{11}=2x_2x_3x_4,\quad H_{12}=2x_1x_3x_4,\quad
H_{13}=2x_1x_2x_4,\quad H_{14}=2x_1x_2x_3,\\
H_{23}=x_1^2x_4,\quad H_{24}=x_1^2x_3,\quad
H_{34}=x_1^2x_2;
\end{gathered}
\]
the other three are zero. The seven displayed monomials are distinct, so
\(u=7\).

For \(s=2\), \(H_{33}=H_{44}=0\), and
\[
H_{11}-H_{12}+H_{22}=0.
\]
The span of \(H_{11},H_{12},H_{22}\) is two-dimensional and supported on
the two distinct squarefree monomials \(x_2x_3x_4,x_1x_3x_4\).
The remaining five nonzero derivatives are
\[
\begin{aligned}
H_{13}&=x_2x_4(2x_1+x_2),&
H_{14}&=x_2x_3(2x_1+x_2),\\
H_{23}&=x_1x_4(x_1+2x_2),&
H_{24}&=x_1x_3(x_1+2x_2),\\
H_{34}&=x_1x_2(x_1+x_2).
\end{aligned}
\]
Their respective square monomials \(x_2^2x_4,x_2^2x_3,
x_1^2x_4,x_1^2x_3,x_1^2x_2\) show that these five forms
are independent and independent of the preceding squarefree span.
Thus \(u=7\).

For \(s=3\), \(H_{44}=0\). The six off-diagonal derivatives are
\[
\begin{aligned}
H_{12}&=x_3x_4(2x_1+2x_2+x_3),&
H_{13}&=x_2x_4(2x_1+x_2+2x_3),\\
H_{14}&=x_2x_3(2x_1+x_2+x_3),&
H_{23}&=x_1x_4(x_1+2x_2+2x_3),\\
H_{24}&=x_1x_3(x_1+2x_2+x_3),&
H_{34}&=x_1x_2(x_1+x_2+2x_3).
\end{aligned}
\]
They are independent: each contains a square monomial not occurring in
any other displayed derivative, namely
\(x_3^2x_4,x_2^2x_4,x_2^2x_3,x_1^2x_4,x_1^2x_3,x_1^2x_2\),
respectively. The remaining three nonzero diagonal derivatives
\(H_{11}=2x_2x_3x_4\), \(H_{22}=2x_1x_3x_4\), and
\(H_{33}=2x_1x_2x_4\) are independent squarefree monomials, so \(u=9\).

For \(s=4\), all six off-diagonal derivatives have the form
\[
H_{ij}=x_kx_l(2x_i+2x_j+x_k+x_l),\qquad
\{k,l\}=[4]\setminus\{i,j\}.
\]
The square monomials \(x_k^2x_l\) and \(x_kx_l^2\) occur only in
the derivative indexed by the complementary pair \(\{i,j\}\).
Consequently the six off-diagonal derivatives are independent.
The four diagonal derivatives are the four distinct squarefree monomials
\(2\prod_{j\ne i}x_j\), hence are independent of one another and of the
off-diagonal derivatives. Thus \(u=10\).
We have proved the exact values \(u=7,7,9,10\) by explicit identities
and coefficient comparison.

For \(s=1,2\), the actual cubic space contains \(x_1x_3x_4\) and
\(x_2x_3x_4\). Every nonzero direction differentiates at least one of
these to a nonzero quadric on at most three variables. That quadric has
rank at most three and hence cannot belong to \(R\). Thus \(\rank\beta\ge1\),
which with \(u=7\) is more than enough.

For \(s=3\), the actual space contains
\(x_1x_2x_4,x_1x_3x_4,x_2x_3x_4\). Write the direction as
\(\lambda=\sum_{j=1}^4a_j\partial_{x_j}\). If \(a_4\ne0\), projecting
their derivatives modulo \(x_4L\) gives three independent pair monomials,
so quotienting by \(R\) leaves rank at least two. If \(a_4=0\), their
derivatives span a space of dimension at least two divisible by \(x_4\):
if, for example, \(a_1\ne0\), the first two have independent linear
factors. Every quadric in this space has rank at most two, so its
intersection with \(R\) is zero. Again \(\rank\beta\ge2\).

For \(s=4\), the actual space contains all four squarefree cubics. Choose
\(a_j\ne0\); differentiating the three cubics containing \(x_j\), modulo
\(x_jL\), gives three independent pair monomials. The quotient by \(R\)
has rank at least two. The stated inequality follows for \(u=9,10\).
All cubic-space inclusions used here follow by direct differentiation.
\end{proof}

\begin{proposition}[Half-defect estimate]\label{prop:local}
For every nonzero quintic Chow term and every quotient of its actual factor
span,
\[\boxed{\rank\beta+\frac{10-u}{2}\ge2d.}\]
\end{proposition}
\begin{proof}
The independent case is Lemma~\ref{lem:independent}. For factor-span
dimension \(\ell\le4\), the middle rank is bounded above by
\(\min\{10,\binom{\ell+1}{2}\}\), because \(u=\dim F\).
It is bounded below by \(1,2,4,7\), respectively for
\(\ell=1,2,3,4\). Indeed, choose \(\ell\) independent factors as
coordinates and specialize every remaining factor to a coordinate initial
by a generic diagonal one-parameter subgroup. The resulting monomial uses
all \(\ell\) variables and has total degree five. Counting its degree-two
divisors, over the positive partitions of five, gives these minima.
Catalecticant rank cannot increase under specialization.

If \(\ell\le3\), every \(d>0\) has complementary span at most two.
Lemma~\ref{lem:kernel}, minimized over the indicated integer interval for
\(u\), gives the following lower bounds for
\(\rank\beta+(10-u)/2\):
\[
\begin{array}{c|c}
\ell& d=0,\ldots,\ell\\\hline
1&(9/2,11/2)\\
2&(7/2,5,6)\\
3&(2,3,6,7).
\end{array}
\]
For \(\ell=4\), the same lemma covers \(d=2,3,4\) and gives
\(9/2,15/2,17/2\), respectively. The case \(d=0\) is immediate, and
\(d=1\) is Lemma~\ref{lem:four}. Every bound is at least \(2d\).
\end{proof}

\section{Two global estimates}
\begin{lemma}[Rectangular image-span inequality]\label{lem:rectangular}
Let \(A_i:X^*\to Y\) be linear maps of ranks \(u_i\). If
\(D_Y=\dim\sum_i\im A_i\) and
\(D_X=\dim\sum_i\im A_i^*\), then
\[\rank\Bigl(\sum_iA_i\Bigr)\ge D_Y+D_X-\sum_i u_i.\]
\end{lemma}
\begin{proof}
Factor \(A_i=B_iC_i\), with \(B_i:\K^{u_i}\to Y\) injective and
\(C_i:X^*\to\K^{u_i}\) surjective. The horizontal concatenation \(B\)
has rank \(D_Y\), and the vertical concatenation \(C\) has rank
\(D_X\), since its dual image is \(\sum_i\im A_i^*\).
Now \(\sum_iA_i=BC\); Sylvester's rank inequality proves the claim.
\end{proof}

Suppose \(f=\sum_{i=1}^N T_i\), with nonzero quintic Chow terms. Put
\[
 F_i=\cD_2(T_i),\quad U_i=\cD_3(T_i),\quad u_i=\dim F_i=\dim U_i,
 \quad \Delta=\sum_i(10-u_i),
\]
\[
 F=\sum_iF_i,\quad U=\sum_iU_i,\quad
 \dim F=100+h_2,\quad \dim U=100+h_3.
\]
Linearity of the catalecticants gives \(E_2\subset F\) and
\(E_3\subset U\), so \(h_2,h_3\ge0\). Under the differential pairings,
the transpose of \(C_{3,2}(T_i)\) is \(C_{2,3}(T_i)\).
Lemma~\ref{lem:rectangular}, applied to their sum, gives
\begin{equation}\label{eq:upper}
 h_2+h_3\le10N-100-\Delta.
\end{equation}

Let \(q:F\to F/E_2\). The map
\[
 \widetilde\beta:\bigoplus_iU_i\longrightarrow V\otimes(F/E_2),
 \qquad (g_i)\longmapsto\sum_i(1_V\otimes q)\delta(g_i)
\]
factors through addition into \(U\). By Lemma~\ref{lem:permanent}, its
kernel after addition is exactly \(E_3\), hence
\(\rank\widetilde\beta=h_3\).

Set \(L_i=\Span\{\text{factors of }T_i\}\), choose any order, and put
\(W_i=L_1+\cdots+L_i\), \(W_0=0\), and
\(d_i=\dim(W_i/W_{i-1})\). These spaces span \(V\), since
\(V=\partial E_2\subset\partial F\subset\sum_iL_i\), so
\(\sum_i d_i=25\). Let \(Z_i\) be the image of the \(i\)-th summand
of \(\widetilde\beta\), and \(H_i=Z_1+\cdots+Z_i\). Projection to
\((W_i/W_{i-1})\otimes(F/E_2)\) kills \(H_{i-1}\). Thus
\[
 \dim H_i-\dim H_{i-1}\ge
 \rank\bigl((P_i\otimes q|_{F_i})\delta|_{U_i}\bigr),
\]
where \(P_i:L_i\to L_i/(L_i\cap W_{i-1})\) has rank \(d_i\).
The injection \(F_i/(F_i\cap E_2)\hookrightarrow F/E_2\) shows that
this is exactly the local symbol from Proposition~\ref{prop:local}.
Summing gives
\begin{equation}\label{eq:lower}
 h_3\ge\sum_i\left(2d_i-\frac{10-u_i}{2}\right)
 =50-\frac\Delta2.
\end{equation}
Combining \eqref{eq:upper} and \eqref{eq:lower} yields
\[h_2+\frac\Delta2\le10N-150.\]
Consequently \(N\ge15\). If \(N=15\), both \(h_2\) and \(\Delta\)
are zero. Then \(F=E_2\), while each \(F_i\) has dimension ten and is
contained in \(E_2\cap\Sym^2L_i\), contradicting the one-dimensional
bound of Lemma~\ref{lem:permanent}. Hence \(N\ge16\).

\section{Upper bound, scope, and exact verification}
Glynn's identity \cite{Glynn2010} is
\[
 \perm_5=\frac1{16}
 \sum_{\substack{\epsilon\in\{\pm1\}^5\\\epsilon_1=1}}
 \left(\prod_{r=1}^5\epsilon_r\right)
 \prod_{c=1}^5\left(\sum_{r=1}^5\epsilon_rx_{rc}\right).
\]
To verify it, expand by row choices in each column. Summation over the four
free signs kills a monomial unless each of rows two through five occurs an
odd number of times; total degree five then forces every row to occur once.
Its coefficient is sixteen before division. This gives a sixteen-term
decomposition. Over an arbitrary characteristic-zero field, any hypothetical
shorter decomposition remains one after extension to its algebraic closure,
where the lower bound applies. The displayed upper bound is rational.
This proves Theorem~\ref{thm:main} over every characteristic-zero field.

The theorem concerns ordinary rank, and proves no border-rank equality.
The ancillary standard-library verifiers remain optional regression tests;
their outputs are not premises of this proof. The 352-case scan is replaced
by the closed counting formula in Lemma~\ref{lem:independent}; the four
normal-form ranks are replaced by coefficient identities in Lemma~\ref{lem:four};
and the partition bounds are proved in the appendix. From the extracted
submission root runFrom the extracted submission root run
\begin{verbatim}
python -B anc/verify_submission.py
\end{verbatim}
The independent finite implementations do not replace the written
specialization, rank-four kernel, image-span, or filtration arguments.
The computational files of the earlier September proof package are
preserved byte for byte under \path{anc/original/}, together with its complete
unchanged manifest. Historical manuscript files are retained in the public
repository rather than in this ancillary directory; the ancillary README
specifies the exact computational subset. The present lower bound uses the
short argument above and does not depend on that package's state routing,
terminal orbit classification, or historical private replay receipts.

\appendix
\section{Coefficient proofs for the local bounds}
For \(T_s=x_1x_2x_3x_4(x_1+\cdots+x_s)\), the exact second-derivative
ranks used in Lemma~\ref{lem:four} follow by coefficient comparison.
For \(s=1\), the seven nonzero derivatives are distinct monomials. For
\(s=2\), two derivatives vanish, the only relation among the remaining
eight is \(H_{11}-H_{12}+H_{22}=0\), and the other five derivatives have
distinct square monomials. For \(s=3\), the six off-diagonal derivatives
have distinct square-monomial pivots, while the three diagonal derivatives
are distinct squarefree monomials. For \(s=4\), the six off-diagonal
derivatives have distinct square-monomial pivots and the four diagonal
derivatives are distinct squarefree monomials. This proves the ranks
\(7,7,9,10\) without a matrix-rank computation.

The lower bounds \(1,2,4,7\) for the cubic derivative-space dimensions
when the factor span has dimensions \(1,2,3,4\) follow by specializing
the five factors to a monomial with a positive exponent partition of five.
The minimum numbers of distinct degree-two divisors occur at partitions
\((5),(4,1),(3,1,1),(2,1,1,1)\), and are respectively \(1,2,4,7\).
Catalecticant rank cannot increase under specialization, so these are
valid lower bounds for the original product.



\section*{Disclosure of AI-assisted tools}
ChatGPT was used to help audit and develop arguments, revise the exposition,
and prepare verification and submission materials. Responsibility for the
final manuscript rests with the author.

\begin{thebibliography}{9}
\bibitem{HanSong2026}
J.~I. Han and J.~Song, \emph{The rank of the \(5\times5\) permanent
tensor is sixteen}, preprint, 2026,
\href{https://arxiv.org/abs/2608.09708}{arXiv:2608.09708}.
\bibitem{Glynn2010}
D.~G. Glynn, \emph{The permanent of a square matrix},
European J. Combin. \textbf{31} (2010), no.~7, 1887--1891,
\href{https://doi.org/10.1016/j.ejc.2010.01.010}{doi:10.1016/j.ejc.2010.01.010}.
\bibitem{IltenTeitler2016}
N.~Ilten and Z.~Teitler,
\emph{Product ranks of the \(3\times3\) determinant and permanent},
Canad. Math. Bull. \textbf{59} (2016), no.~2, 311--319,
\href{https://doi.org/10.4153/CMB-2015-076-1}{doi:10.4153/CMB-2015-076-1}.
\end{thebibliography}
\end{document}
